% GENERATED FILE. Edit canonical lessons, then run npm run generate:books.
% canonical-source-hash: fnv1a64:ed0cd71935ed4ad6
% canonical-lessons: GE-C122-knoblauch, GE-C122-mehl, GE-C122-schokolade, GE-C122-keks, GE-C122-broetchen

\chapter{Garlic, Flour, Chocolate, Biscuit, Bread roll}
\label{ch:ge-garlic-flour-chocolate-biscuit-bread-roll}
\hlchaptermodality{122}

\begin{chapteropening}
\textbf{By the end of this chapter:} \emph{I can say garlic, flour, chocolate, a biscuit and a bread roll.}

Complete the last lesson of chapter 122: i can say garlic, flour, chocolate, a biscuit and a bread roll.
\end{chapteropening}

\section[der Knoblauch]{der Knoblauch — garlic}

\label{lesson:GE-C122-knoblauch}



\begin{quote}
Before the new one: say the German for a cherry, then the German for a melon.
\end{quote}

\subsection*{You'll want to know: der Knoblauch}
\textbf{der Knoblauch} — \textquotedblleft{}garlic\textquotedblright{}.

\begin{grammarlens}[title={What you've built}]
One more word for this chapter.
\end{grammarlens}

\subsection*{Your turn}
\begin{itemize}
\raggedright
  \item \emph{Say it:} \emph{der Knoblauch}
  \item \emph{Say it:} \emph{der Knoblauch}, once more
  \item \emph{Say it:} say \emph{der Knoblauch}
  \item \emph{Recall:} say the German for a cherry, then the German for a melon, then say \emph{der Knoblauch} again
\end{itemize}

\begin{tcolorbox}[breakable,colback=teal!4,colframe=teal!35!black,title={Before you move on}]
What does der Knoblauch mean? (\textquotedblleft{}Garlic\textquotedblright{}.) Say it once more.
\end{tcolorbox}
\section[das Mehl]{das Mehl — flour}

\label{lesson:GE-C122-mehl}



\begin{quote}
Before the new one: say the German for a melon, then the German for garlic.
\end{quote}

\subsection*{You'll want to know: das Mehl}
\textbf{das Mehl} — \textquotedblleft{}flour\textquotedblright{}.

\begin{grammarlens}[title={What you've built}]
One more word for this chapter.
\end{grammarlens}

\subsection*{Your turn}
\begin{itemize}
\raggedright
  \item \emph{Say it:} \emph{das Mehl}
  \item \emph{Say it:} \emph{das Mehl}, once more
  \item \emph{Say it:} say \emph{das Mehl}
  \item \emph{Recall:} say the German for a melon, then the German for garlic, then say \emph{das Mehl} again
\end{itemize}

\begin{tcolorbox}[breakable,colback=teal!4,colframe=teal!35!black,title={Before you move on}]
What does das Mehl mean? (\textquotedblleft{}Flour\textquotedblright{}.) Say it once more.
\end{tcolorbox}
\section[die Schokolade]{die Schokolade — chocolate}

\label{lesson:GE-C122-schokolade}



\begin{quote}
Before the new one: say the German for garlic, then the German for flour.
\end{quote}

\subsection*{You'll want to know: die Schokolade}
\textbf{die Schokolade} — \textquotedblleft{}chocolate\textquotedblright{}.

\begin{grammarlens}[title={What you've built}]
One more word for this chapter.
\end{grammarlens}

\subsection*{Your turn}
\begin{itemize}
\raggedright
  \item \emph{Say it:} \emph{die Schokolade}
  \item \emph{Say it:} \emph{die Schokolade}, once more
  \item \emph{Say it:} say \emph{die Schokolade}
  \item \emph{Recall:} say the German for garlic, then the German for flour, then say \emph{die Schokolade} again
\end{itemize}

\begin{tcolorbox}[breakable,colback=teal!4,colframe=teal!35!black,title={Before you move on}]
What does die Schokolade mean? (\textquotedblleft{}Chocolate\textquotedblright{}.) Say it once more.
\end{tcolorbox}
\section[der Keks]{der Keks — a biscuit}

\label{lesson:GE-C122-keks}



\begin{quote}
Before the new one: say the German for flour, then the German for chocolate.
\end{quote}

\subsection*{You'll want to know: der Keks}
\textbf{der Keks} — \textquotedblleft{}a biscuit\textquotedblright{}.

\begin{grammarlens}[title={What you've built}]
One more word for this chapter.
\end{grammarlens}

\subsection*{Your turn}
\begin{itemize}
\raggedright
  \item \emph{Say it:} \emph{der Keks}
  \item \emph{Say it:} \emph{der Keks}, once more
  \item \emph{Say it:} say \emph{der Keks}
  \item \emph{Recall:} say the German for flour, then the German for chocolate, then say \emph{der Keks} again
\end{itemize}

\begin{tcolorbox}[breakable,colback=teal!4,colframe=teal!35!black,title={Before you move on}]
What does der Keks mean? (\textquotedblleft{}A biscuit\textquotedblright{}.) Say it once more.
\end{tcolorbox}
\section[das Brötchen]{das Brötchen — a bread roll}

\label{lesson:GE-C122-broetchen}



\begin{quote}
Before the new one: say the German for chocolate, then the German for a biscuit.
\end{quote}

\subsection*{You'll want to know: das Brötchen}
\textbf{das Brötchen} — \textquotedblleft{}a bread roll\textquotedblright{}.

\begin{grammarlens}[title={What you've built}]
One more word for this chapter.
\end{grammarlens}

\subsection*{Your turn}
\begin{itemize}
\raggedright
  \item \emph{Say it:} \emph{das Brötchen}
  \item \emph{Say it:} \emph{das Brötchen}, once more
  \item \emph{Say it:} say \emph{das Brötchen}
  \item \emph{Recall:} say the German for chocolate, then the German for a biscuit, then say \emph{das Brötchen} again
\end{itemize}

\begin{tcolorbox}[breakable,colback=teal!4,colframe=teal!35!black,title={Before you move on}]
What does das Brötchen mean? (\textquotedblleft{}A bread roll\textquotedblright{}.) Say it once more.
\end{tcolorbox}
